%% kapa: exam=1 round=23 year=2012 period=2 form=A subject=appraisal_law subject_ko=부동산관계법규 seq=11 number=51
%% source: Q-Net 감정평가사 기출문제 (cst00309) artlSeq=5052580 file=raw/1008/5052580/2012년 제23회 감정평가사 제1차시험 기출문제(2교시 A형).pdf
%% conversion: kordoc 4.24.51 -> tools/gen1.py (auto); review=proofread-sonnet; flags=none
\begin{question}{51}
국토의 계획 및 이용에 관한 법령상 토지거래허가 구역에 관한 설명으로 옳은 것은?
\begin{choices}
  \item 허가구역의 지정권자는 국토해양부장관이지만 특별히 투기가 성행할 우려가 있을 경우에는 관계 중앙행정기관의 장도 지정할 수 있다.
  \item 국토해양부장관이 허가구역을 지정하려면 중앙도시계획위원회의 심의를 거쳐야 한다.
  \item 허가구역을 처음으로 지정하고자 할 때에는 미리 시ㆍ도지사 및 시장ㆍ군수 또는 구청장의 의견을 들어야 한다.
  \item 허가구역의 지정은 허가구역의 지정을 공고한 날부터 그 효력이 발생한다.
  \item 허가구역의 지정에 대해서는 15일 이상 그 사실을 공고하고 공고내용을 15일 이상 일반이 열람할 수 있도록 하여야 한다.
\end{choices}
\end{question}
\answer{2}
